\chapter[Impacts of BSSs]{Impacts of \acp{bss}}
\label{ch:impacts_of_bss}

\acp{bss} have attracted significant scholarly and policy attention not only for their role as an innovative mobility service but also for their potential impacts across multiple dimensions of urban life. Their rapid growth has been driven by promises of diverse benefits: improved urban mobility, environmental sustainability, public health gains, and new economic opportunities. At the same time, questions have been raised about equity of access, safety, and the overall magnitude of their impacts.

\textcolor{blue}{
    The contributions in Parts~\ref{pt:data-driven-station-planning} and~\ref{pt:predictive_operations} develop methods for station planning and fleet operations rather than study and quantify these impacts; what follows nonetheless clarifies how \acp{bss} interventions are usually evaluated, so that the contributions of this thesis can be read alongside the expectations established in the literature.
}
Accordingly, this Chapter synthesizes the research on the global impacts of \acp{bss} over the past 15 years, focusing on mobility, environmental, health, and socioeconomic dimensions.


\section{Mobility and transport impacts}
\label{sec:mobility_and_transport_impacts}

This Section reviews the main mobility impacts of \acp{bss}, focusing on how they shape trip patterns, interact with \ac{pt}, alter modal share, and influence broader transport system performance.

    \subsection{Trip patterns and multimodal integration}

Bike-sharing has measurably altered urban mobility patterns by adding a flexible, short-trip transport option that often complements existing modes. Evidence consistently shows that \acp{bss} are used predominantly for utilitarian purposes, with commuting and errands being the most common, while also supporting leisure and tourism mobility~\cite{Ricci2015}.

% Time efficiency gains
A frequently cited benefit of this utilitarian usage is improved first- and last-mile connectivity to \ac{pt}~\cite{Shaheen2010, Qiu2018}. By strategically siting bike stations near \ac{pt} hubs, \ac{bss} can extend the reach of metro and bus networks, and therefore facilitate multimodal trips~\cite{Cheng2022}. In Beijing, for example, bike-sharing access to bus and metro stations saved commuters an average of eight minutes per day, translating into substantial citywide time savings and even modest productivity gains~\cite{Qiu2018}. Recent analyses also point to additional efficiency gains from \acp{e-b}, whose travel speeds are approximately 9\% higher than those of \acp{m-b}, which is equivalent to roughly 70 seconds saved on a typical 1.5–2 km trip~\cite{Waldner2025}, further highlighting the contribution of contemporary \acp{bss} to urban time efficiency. 

    \subsection{Modal substitution and car-use reduction}

% Modal substitution
Beyond connectivity and time efficiency, \acp{bss} capacity for modal substitution is central, as it is one of the most debated aspects of their impact. Understanding which modes bike-sharing replaces is essential, as reducing car dependency is a core motivation for their introduction~\cite{Fishman2016}. Early surveys of the 2010s across cities including Dublin, London, Lyon, Barcelona, Montreal, Washington D.C., Minneapolis, Melbourne and Brisbane consistently showed that only 5--20\% of bike-sharing trips would otherwise have been made by car, with the majority substituting for walking or \ac{pt}~\cite{Murphy2015, Fishman2013, Fishman2014, BachandMarleau2012}. Ricci~\cite{Ricci2015} concluded that there was “no evidence that bike-sharing significantly reduces traffic congestion, carbon emissions and pollution”. At the same time, she cautioned that examining mode substitution alone does not capture the full range of impacts. To obtain a more complete assessment, she argued that research should also consider the characteristics of replaced car trips (e.g., frequency, duration, route) and the extent of induced demand through new trips that would not otherwise have been made—factors that were rarely examined in the earlier survey-based literature. On this basis, both Ricci and Fishman called for more detailed research on this issue in 2015~\cite{Ricci2015, Fishman2016}.  

Subsequent research broadly supports the finding that only a minority of bike-sharing trips replace car travel. Despite the technological advances of \ac{e-b} and dockless systems, the consensus remains the same: while bike-sharing does replace some car usage, the majority of trips continue to come from walking and \ac{pt}. In Beijing a survey of dockless system users found that only 14.8\% of trips replaced car travel~\cite{Chen2022b}; in Delft, the introduction of both docked (\textit{OV-fiets}) and dockless (\textit{Mobike}) systems showed similar patterns, with approximately 30\% of trips replacing car use~\cite{Ma2020}; and a multi-city European survey in 2021--2022 placed the figure at about 21\%~\cite{Bosehans2024}. Although these figures are higher than early estimates, they still indicate that most bike-sharing usage comes at the expense of active modes or \ac{pt} rather than driving, and that some authors argue that \ac{e-b} are not enough to attract car users to \acp{bss}~\cite{Teixeira2023}. Accordingly, some authors argue that more restrictive measures, such as car-free zones or increased parking charges, are necessary to strengthen the car-reducing effects of \acp{bss}~\cite{Bosehans2024}.

Even when bike-sharing does not fully replace car journeys, it can reduce car dependence indirectly by enabling multimodal travel. Many users combine bike-sharing with \ac{pt} for first- and last-mile segments, and repeated participation in bike-sharing has been associated with reductions in driving frequency. In Washington D.C., for example, 55\% of respondents to the \textit{CaBi} 2016 annual survey reported driving less often since joining the system~\cite{NACTO2020}.

    \subsection{Contextual factors shaping modal substitution}

% Modal substitution context matters
Context, however, matters. Research highlights that \ac{bss} user demographics and trip purposes shape modal substitution outcomes. Dockless systems in China and Europe often attract younger riders and those without cars or driving licenses, meaning these users were unlikely to be replacing car trips in the first place~\cite{Ma2020, Bosehans2024}. By contrast, in cities or niches where bike-sharing directly competes with car convenience, higher car substitution rates are possible. Moreover, cargo bike-sharing services--which allow users to haul goods or children--show particularly strong potential to displace car use: one survey of cargo bike-sharing users found that 46\% would have driven a car if the cargo bike had not been available~\cite{Becker2018}. This suggests that tailoring bike-sharing offerings (e.g. \acp{e-b} for longer commutes, cargo bikes for errands) can increase the likelihood of replacing car trips.

    \subsection{System-level effects of modal substitution}

% Congestion, parking and car ownership reduction
Although direct car substitution rates remain modest, even small shifts can yield measurable system-level benefits. In Washington D.C., \textit{Capital Bike-sharing} reduced neighborhood traffic congestion by about 4\%, with the greatest reductions in already congested areas~\cite{Hamilton2018}. A study of 96 U.S. cities similarly found that bike-sharing deployment eased congestion in large, transit-rich cities, particularly at rush hours, though effects were weaker or mixed in smaller and high-car-ownership contexts~\cite{Wang2017}. In China’s megacities such as Beijing, Shanghai, and Wuhan, the initial rollout of dockless \acp{bss} significantly reduced congestion in the short term, though subsequent expansions had diminishing returns~\cite{Xu2024}. Notably, in these Chinese cities, metro ridership rose alongside bike-sharing use, suggesting that many users shifted from driving to bike+metro combinations. Parking impacts follow a similar logic: in Pittsburgh, the opening of a \ac{bss} station was associated with a 2\% decline in parking meter use~\cite{Pelechrinis2016}, while in Boston the installation of new stations led to a 2.2\% reduction in household vehicle ownership and a 3.3\% drop in vehicle-miles traveled, with the strongest impacts (around 10\% lower car dependence) observed near transit stations~\cite{Basu2021}. Together, these findings suggest that \acp{bss}, as part of a multimodal system, can support “car-lite” lifestyles and free up curb space over the long term.  

% Induced demand
Beyond substitution, \acp{bss} also induce new trips that would not otherwise occur. A European study found that about 11\% of shared mobility trips were entirely new journeys induced by the service~\cite{Bosehans2024}. Such trips do not reduce any particular mode, but they indicate that bike-sharing provides additional mobility and access, expanding travel opportunities for its users. 

    \subsection{Transport network resilience}

% Transport network resilience
Finally, another key mobility impact of \acp{bss} is their contribution to the transport network resilience. During four London metro strikes, bike-sharing trips spiked sharply, with many people who would have used the metro switching to shared bikes, often using longer or unfamiliar routes when their regular stations were inaccessible~\cite{Yang2022}. Similarly, during a major flood in Porto Alegre, Brazil, the local \ac{bss} continued to function and enabled many daily activities to proceed despite widespread transport disruption~\cite{Araujo2025}. Evidence from Washington D.C. further shows that bike-sharing usage rises significantly during planned Metro shutdowns, demonstrating its role as an alternative mode in complete transit outages~\cite{Cheng2022}. Last, the COVID-19 pandemic further reinforced this role of \acp{bss} as a transport mode that brings resilience to the transport network. While \ac{pt} usage remained far below pre-pandemic levels, bike-sharing volumes nearly returned to normal once restrictions eased, with some even increasing its trips numbers~\cite{Teixeira2024}. For instance, in Barcelona's \textit{Bicing}, trip levels exceeded pre-pandemic counts within just a couple of months of reopening after the quarantine~\cite{Grau-Escolano2024}.
\section{Environmental impacts}
\label{sec:environmental_impacts}

The environmental impacts of \acp{bss} have been a central part of their justification~\cite{Shaheen2010, Fishman2016, Midgley2011}. These impacts extend beyond avoided transport emissions and include energy use, material consumption, logistics, maintenance, and end-of-life waste and recycling. While early academic discussions often focused on \ac{ghg} reductions from substituting car trips, recent research has broadened the perspective to include the full life cycle of shared bicycles, docking infrastructure, batteries, and logistics operations.

    \subsection{Avoided transport emissions through modal substitution} 

    Transport emissions comprise both \acp{ghg} (e.g., \ce{CO2}, \ce{CH4}, \ce{N2O}), which accumulate in the atmosphere and drive climate change, and local air pollutants (e.g., \ce{NOx}, \ce{HC}, \ce{PM_{2.5}}, \ce{PM_{10}}, \ce{CO}), which are short-lived and affect urban air quality. Early studies typically estimated avoided emissions by combining (i) cycling distance, (ii) a counterfactual mode that would otherwise have been used, and (iii) mode-specific emission factors, with cycling treated as negligible at the tailpipe. When this approach was combined with simple distance rules for modal substitution and strong extrapolation to the city or fleet scale, it produced very large headline \ce{CO2} savings in systems such as Barcelona's \textit{Bicing}~\cite{Rojas-Rueda2011}, Shanghai's \textit{Mobike}~\cite{Zhang2018}, Beijing's public system~\cite{Qiu2018}, and New York's \textit{Citi Bike}~\cite{Chen2022c}. The savings magnitude is sensitive to how distance is measured (e.g., straight-line vs.\ network routing), but it is mainly driven by substitution assumptions: multi-city surveys typically find car replacement on the order of 5--20\% of trips~\cite{Murphy2015, Fishman2013, Fishman2014}, whereas some published scenarios assumed much higher shares, including 75\%~\cite{Qiu2018} and 90\%~\cite{Rojas-Rueda2011}. If walking or \ac{pt} would have been the true counterfactual, assigning a car counterfactual instead inflates avoided \ce{CO2}, because the emissions difference between cycling and walking or \ac{pt} is much smaller than between cycling and private motorized travel.

    Accordingly, subsequent studies have sought to ground substitution in observed behavior and travel surveys rather than fixed rules, in line with the evidence reviewed in Section~\ref{sec:mobility_and_transport_impacts}. Kou et al.~\cite{Kou2020} developed a trip-level probabilistic assignment that combines observed bike-sharing distances with city-specific travel surveys and life-cycle emission factors, assigning each trip to a most likely replaced mode; applied to eight U.S. cities in 2016, it yields far smaller annual savings than many previous assessments. Saltykova~\cite{Saltykova2022} used traveler surveys for a dockless system in Chengdu, China, finding that 39\% of trips replaced bus, 13.5\% replaced subway, and only 10.9\% replaced car trips. With these figures, avoided emissions fell markedly compared to a ``cars-only'' counterfactual, illustrating how omitting \ac{pt} can inflate climate benefits by a factor of two or more. 
    
    A more refined approach was proposed by Li et al.~\cite{Liu2021}, who developed a high-resolution discrete-choice framework matching over 23 million dockless bike-sharing trips in Shanghai to realistic modal alternatives using travel time, distance, and cost information from online routing services. Their results showed that most trips substituted for walking (30.7\%) or \ac{pt} (50\%), with car replacement below 2\%. This progression from simple distance rules to survey-informed probabilities, and finally to discrete-choice models, illustrates a gradual shift toward more behaviorally grounded and context-sensitive estimates of environmental benefits.
    
    Research has also estimated reductions in local air pollutants (e.g., \ce{NOx}, \ce{SO2}, particulate matter) following the same logic as with \ac{ghg} emissions, though such studies remain relatively uncommon and inherit the same substitution caveats as \acp{ghg} accounts~\cite{Zhang2018, Qiu2018}.



    \subsection[LCA for BSSs]{\ac{lca} for \acp{bss}}

    Avoided transport emissions capture only part of the environmental picture, and \acp{bss} operations can even offset expected benefits: Fishman, Washington, and Haworth~\cite{Fishman2014} showed that London’s \ac{bss} in 2012 increased overall motor vehicle usage due to rebalancing. Consequently, recent work relies on \acp{lca}--typically aligned with ISO 14,040 and 14,044~\cite{LCA2006a, LCA2006b}--to account for burdens from manufacturing, operation, maintenance, and end-of-life management, while also incorporating modal substitution within a single framework (Figure~\ref{fig:LCA_bss}). Results are usually reported per service delivered as g~\ce{CO2}-eq/pkm~\cite{Luo2019, Chen2024, Bonilla-Alicea2020}, where \ce{CO2}-eq aggregates the warming effect of multiple \acp{ghg} (e.g., \ce{CO2}, \ce{CH4}, \ce{N2O}) into an equivalent amount of \ce{CO2} based on global warming potentials. Reported values for complete \ac{bss} \acp{lca} span roughly 30--150~g~\ce{CO2}-eq/pkm~\cite{Zhu2023, Felipe-Falgas2022, Cazzola2020, Sun2022, Calan2024}, reflecting differences in system type, utilization and operations, and included life-cycle phases.

    \begin{figure}[htbp!]
        \centering
        \includegraphics[width=1\textwidth, trim={2cm 1cm 2.5cm 1.5cm}, clip]{0_plots/figure_LCA.pdf}
        \caption[LCA components for BSSs]{
            \textbf{\ac{lca} components for \acp{bss}.} 
            This figure summarises the main stages and impact categories considered in an \ac{lca}. Based on Luo et al.~\cite{Luo2019}, adapted and elaborated by the author.
        }
        \label{fig:LCA_bss}
    \end{figure}

    Manufacturing includes raw material extraction and production of bicycles, docking stations, and (when applicable) batteries. Manufacturing is often the largest single contributor in docked systems, accounting for around 50\% of total life-cycle emissions~\cite{Calan2024, Sun2022}. For \acp{e-b}, hardware production represents 61\% of life-cycle emissions in docked systems and 52\% in dockless systems~\cite{Luo2019}, with batteries contributing 20--30\% of embodied \acp{ghg}~\cite{Zhu2023}. Material choices also shape embodied impacts: for private bicycles, aluminum frames have been reported to carry higher production impacts than alternatives such as steel~\cite{Dave2010, Leunberger2010, Luna2016} or bamboo~\cite{Agyekum2017}.

    While cycling itself is practically emission-free, the use phase includes supporting operations, with rebalancing consistently emerging as a dominant contributor. Luo et al.~\cite{Luo2019} estimated that rebalancing accounts for 36\% of total life-cycle emissions in docked systems and up to 73\% in dockless fleets~\cite{Luo2019}. Impacts vary with fleet size, system scale, vehicle type, and redistribution frequency~\cite{Bonilla-Alicea2020}, with reported service-vehicle intensities ranging from 0.012~km of van travel per bike-km in docked systems~\cite{Felipe-Falgas2022} to 0.10--0.26~km in dockless fleets~\cite{Sun2022, Cazzola2020}. Vehicle technologies further shape emissions, from conventional vans~\cite{Chen2020LifeCycle, Hollingsworth2019} to electric vans or cargo bikes for battery collection and charging~\cite{Felipe-Falgas2022, Calan2024}. For electrically assisted systems, charging adds additional impacts: Sun et al.~\cite{Sun2022} reported electricity demands with climate impacts ranging from about 10--20~g~\ce{CO2}-eq/km in low-carbon grids to more than seven times higher in coal-dominated regions~\cite{Zhu2023, Calan2024}.  Maintenance is often underrepresented due to limited operator records~\cite{Calan2024}, but when included it can contribute around 10--15\% of total life-cycle \acp{ghg} in docked systems, with larger shares in dockless fleets~\cite{Sun2022, Chen2024, deBortoli2021Environmental, Luo2019}. In here, assumptions about component lifetimes are central, and improving durability and repairability can mitigate environmental impacts~\cite{Zhu2023, Calan2024}. Most \acp{lca} also incorporate modal substitution by assigning each trip a counterfactual mode and estimating avoided life-cycle emissions; outcomes depend critically on the share of car and ride-hailing trips displaced versus low-carbon modes~\cite{Krauss2022}. For instance, in Barcelona, Felipe-Falgas et al.~\cite{Felipe-Falgas2022} show that shared \acp{e-b} can even increase \acp{ghg} when they mainly replace \ac{pt} or walking rather than car trips.

    End-of-life impacts are often underrepresented due to data limitations~\cite{Calan2024} and typically represent only a small share of total life-cycle emissions~\cite{Krauss2022}. Beyond formal disposal, fleet attrition and oversupply can exacerbate impacts: Calan et al.~\cite{Calan2024} noted that losing around 20\% of bikes annually can increase total fleet emissions by as much as 25\%. A dramatic demonstration of this challenge was seen in Chinese cities during the 2017--2018 dockless bike boom, when ``bike-sharing graveyards'' (Figure~\ref{fig:bss_graveyard}) emerged as excess and abandoned bikes accumulated in massive lots~\cite{Sun2022, AmusingPlanet2018}. In these situations, low utilization prevented manufacturing emissions from being distributed across enough trips, and inadequate recycling meant much of the material value was lost~\cite{Luo2019, Sun2022}.

\begin{figure}    
    \centering
    \includegraphics[width=0.8\textwidth]{0_plots/graveyard.jpg}
    \caption[BSS graveyard in Xiamen, China]{
        \textbf{\ac{bss} graveyard in Xiamen, China.} 
        Photograph sourced from \textit{Amusing Planet}~\cite{AmusingPlanet2018} illustrating large-scale accumulation of shared bikes.
    }
    \label{fig:bss_graveyard}
\end{figure}


Overall, the environmental performance of \acp{bss} depends on system design, operational practices, and local mobility patterns. Benefits from avoided transport emissions materialize primarily when car substitution is substantial and operational burdens are minimized, particularly rebalancing.









\section{Health and safety impacts}
\label{sec:health_and_safety_impacts}

Health is also a critical dimension in \acp{bss} promotion, complementing the environmental and transport outcomes discussed earlier~\cite{DeMaio2009, Fishman2016, Qiu2018}. Regular cycling is consistently associated with improved cardiovascular fitness and lower risks of chronic disease and premature mortality~\cite{Green2021, Celis-Morales2017, Dutheil2020}, and it provides an efficient means of achieving the recommended 150 minutes of moderate to vigorous physical activity per week~\cite{Knowler2002}. Cycling can also support mental well-being~\cite{Chekroud2018, Ruiz2015}, but evidence that bike-sharing substantially increases physical activity at the population level is mixed. Early observational and survey-based studies reported increased cycling among residents near stations and self-reported health improvements among users~\cite{Fuller2013, Molina-Garcia2013, Alberts2012}, and county-level analyses linked \ac{bss} deployment with small declines in obesity prevalence~\cite{Xu2019, Hosford2019}. However, the magnitude of health gains depends strongly on which modes are replaced~\cite{Fishman2014, Murphy2015}. Accordingly, a narrative review concluded that many systems may have limited effects on overall population physical activity because \ac{bss} users are often already active and trips can substitute for walking or \ac{pt} rather than car travel~\cite{Bauman2017}.

To quantify net population-level effects and explicitly account for both benefits and risks, researchers commonly use \acp{hia}, which combine travel behavior data with epidemiological dose--response relationships~\cite{Rojas-Rueda2011, Woodcock2014, Otero2018, Clockston2021}. Here, \acp{hia} typically model three pathways~\cite{Teixeira2021}: (i) health gains from increased physical activity, (ii) health risks from traffic injuries, and (iii) health losses from exposure to air pollution, most often operationalized through in-transit \ce{PM_{2.5}} exposure~\cite{Woodcock2014, Otero2018, Clockston2021, Cao2021}. Results are sensitive to key assumptions, especially modal substitution, and to whether assessments account only for users' in-transit exposure or also include system-wide air-quality changes from reduced car use~\cite{Clockston2021, Hou2023}.

Across empirical applications, \acp{hia} consistently find that net health impacts are driven primarily by gains in physical activity, with air-pollution and traffic-injury pathways generally playing smaller roles. In Barcelona, Rojas-Rueda et al.~\cite{Rojas-Rueda2011} estimated that \textit{Bicing} prevented approximately 12.5 deaths per year, with comparatively small risks from traffic injuries and pollution exposure, yielding a benefit--risk ratio of about 77:1. However, this result partly reflected an optimistic modeling assumption of 90\% of \textit{Bicing} trips replacing car journeys, a limitation they acknowledged. In London, Woodcock et al.~\cite{Woodcock2014} similarly found a positive overall health balance using empirical trip data and disaggregation by age and sex, and reported heterogeneity across demographic groups: men experienced larger benefits because they cycled more frequently and for longer distances, older adults benefited more in absolute terms due to higher baseline mortality risk, and women faced higher relative injury risks in traffic.
 
Extending the evidence across 12 European systems, Otero et al.~\cite{Otero2018} reported a benefit--risk ratio of 19:1 and estimated 5.17 deaths avoided per 1,000 bicycles in service per year, while emphasizing that benefits could increase substantially under scenarios with higher car substitution. In the U.S., Clockston et al.~\cite{Clockston2021} estimated that \acp{bss} use prevented about 4.7 premature deaths annually nationwide (approximately \$36 million in avoided health costs) and around 2 premature deaths annually in New York City (about \$15 million). Beyond Europe and North America, Chen et al.~\cite{Chen2023} evaluated Shanghai’s \textit{Mobike} program in a highly polluted megacity context. For reference, in 2021, the \ac{who} recommended long-term average \ce{PM_{2.5}} concentrations not exceeding 5~\si{\micro g/m^3} and daily averages below 15~\si{\micro g/m^3}~\cite{WHO2021}. Despite average \ce{PM_{2.5}} concentrations of about 45~\si{\micro g/m^3}, their \ac{hia} found that health gains from increased physical activity still outweighed the added pollution burden, preventing roughly 7--8 premature deaths per year among about 306,000 users. The net balance began to diminish only as concentrations approached about 68~\si{\micro g/m^3}~\cite{Chen2023}.

The air-pollution pathway is typically marginal in standard \acp{hia} that focus on users' in-transit exposure~\cite{Rojas-Rueda2011, Woodcock2014, Otero2018, Clockston2021}, though it can become more influential in highly polluted contexts~\cite{Chen2023}. Some work has also emphasized that exposure can vary sharply over space and time, and that daily mean concentrations may obscure short-term peaks relevant to cyclists~\cite{Cao2021}.

Similarly, traffic-injury impacts generally remain small relative to physical-activity benefits, but they vary across cities and are sensitive to local safety environments and infrastructure~\cite{Rojas-Rueda2011, Woodcock2014, Otero2018, Clockston2021}. Evidence from London suggests that injury rates for bike-sharing users can be at least as low as those for private cyclists~\cite{Woodcock2014}. Surveys likewise report lower crash rates for bike-sharing than for private cycling and, in some analyses, lower risks of fatal or severe injuries after scheme introduction~\cite{Martin2016, Fishman2016_road_safety}. Indeed, reviews highlight a ``safety-in-numbers'' effect, whereby increases in cycling volumes can reduce the risk of injury per cyclist, helping to explain why higher ridership due to the existence of a \acp{bss} does not necessarily translate into proportionate increases in injuries~\cite{Jacobsen2015, Teixeira2021}. At the same time, a systematic review reported high rates of riding without a helmet among bike-sharing users and mixed evidence on whether injury rates change after scheme launches~\cite{Chen2021}. One multi-city analysis further found that in jurisdictions without mandatory helmet laws, the share of head-injured cyclists increased after bike-sharing implementation~\cite{Graves2014}. 

Finally, although \acp{bss} deliver clear population-level health benefits, these gains are not necessarily distributed equally across urban space or social groups. In New York City, station placement has historically favored higher-income neighborhoods, shaping who can access and benefit from the system; while network expansion increased estimated premature deaths averted (from two to three per year), sensitivity analyses suggested that prioritizing higher-poverty areas could have yielded larger benefits~\cite{Babagoli2019}. National-level \acp{hia} evidence from the U.S. similarly emphasizes that equity depends on more balanced station distribution and improved data on users’ socioeconomic characteristics~\cite{Clockston2021}. More broadly, because most evidence comes from high-income cities with relatively low air pollution and stronger road-safety environments, the benefit--risk balance may differ in settings with poorer air quality or weaker safety conditions, underscoring the need to expand the evidence base across contexts~\cite{Teixeira2021, Tainio2016}.
\section{Socioeconomic impacts}

\section{Economic impacts}

The economic impacts of \acp{bss} are multi-layered, and while most \acp{bss} operate with public support or sponsorship and typically do not generate direct profits~\cite{Luo2019}, their economic relevance becomes more apparent once externalities such as travel time savings, health improvements, and local spending are considered.

At the user level, bike-sharing can reduce both monetary and time costs, especially for users shifting from car, ride-hailing, or \ac{pt}. However, because many trips replace walking or \ac{pt} rather than car use~\cite{Murphy2015, Fishman2013, Fishman2014, BachandMarleau2012}, direct per-trip financial savings are often limited, whereas time savings can still be meaningful given the economic value typically assigned to travel time savings~\cite{Qiu2018}. Aggregated across users, these effects can be substantial. In Beijing, bike-sharing as a last-mile commuting solution was estimated to save about 8 minutes per worker per day, yielding annual productivity gains of 0.3--1.2 billion yuan depending on assumptions~\cite{Qiu2018}. In Dublin, \textit{Dublinbikes} was associated with annual savings of 16.5 million minutes compared to prior trips (often walking or \ac{pt}), valued at about 6~million~\texteuro{} in direct time benefits, plus an additional 6.79~million~\texteuro{} per year in wider agglomeration benefits through increased accessibility and effective density~\cite{Bullock2017}. Evidence from London similarly indicates \ac{bss} trips are, on average, about 20\% shorter than the journeys they replace~\cite{Woodcock2014}, and U.S. evidence suggests bike-sharing can be as fast or faster than taxis for short \ac{od} trips (under 3~km) during rush-hour congestion~\cite{FaghihImani2017_Hail}. Overall, these findings suggest that bike-sharing can contribute to more time-efficient urban mobility.

Beyond travel-time effects, bike-sharing can stimulate local commerce by improving access and enabling additional trips. Survey evidence from \textit{Nice Ride Minnesota} suggests users made extra trips--primarily to retail and dining destinations--because of bike-sharing convenience, reporting nearly \$35,000 in expenditures in one week~\cite{Schoner2012}. In Washington, D.C., \textit{Capital Bikeshare} users reported making trips they would not otherwise have taken (16\%) and spending more at destinations reached by bike (23\%), while businesses near stations reported increased sales and positive neighborhood impacts~\cite{Buehler2014}. Although some effects may reflect redistribution rather than net new spending, the evidence supports bike-sharing as a contributor to local commercial activity and visibility. Moreover, bike-sharing has also been linked to property-market effects. In Montreal, proximity to stations was associated with higher multifamily housing sale prices. Each additional station within 800 meters corresponded to about a \$709 increase in unit price, implying roughly \$8,650 per unit (about 2.7\% of average prices) given typical station densities in the study area~\cite{ElGeneidy2016}.

Health benefits add a further economic dimension. Across 12 European cities, annual health benefits attributable to \acp{bss} use have been estimated at 18 million~\texteuro{}~\cite{Otero2018}, while in the U.S. annual health-related economic gains have been valued at around \$36 million~\cite{Clockston2021}. A cost--benefit analysis of \textit{Dublinbikes} found that benefits--notably health-related--outstripped costs, with benefit-cost ratios ranging from 1.21 to 1.68, ensuring a net societal gain for every euro invested~\cite{Murphy2015}.

Despite these diverse positive impacts, potential negative economic consequences should not be overlooked. The rapid expansion of dockless \acp{bss} in China around 2017 illustrates such risks: some major operators experienced mass layoffs and withdrew from multiple cities~\cite{Fortune2017, BusinessInsider2018}, and public agencies eventually bore costs related to abandoned bikes and urban clutter~\cite{Chen2019}. Thus, as cities leverage the economic upsides of bike-sharing, prudent planning and regulation remain important to mitigate wider costs.

\section{Social equality}

While \acp{bss} are frequently promoted as instruments for advancing social equity~\cite{Shaheen2010}, their ability to do so depends heavily on who can access and use them. Because proximity to stations strongly predicts usage~\cite{Kabra2018, BachandMarleau2012}, unequal spatial coverage can reinforce existing mobility inequalities, particularly among disadvantaged groups affected by transport poverty (e.g., low car ownership and limited \ac{pt} access)~\cite{Lucas2016}. Accordingly, research consistently identifies spatial distribution as a key driver of inequality in bike-sharing~\cite{Teixeira2021, Goodman2014, Ogilvie2012}. Stations are often concentrated in central, dense, and affluent neighborhoods, leaving peripheral or lower-income areas with weaker access, a pattern documented in cross-city spatial fairness analyses~\cite{DuranRodas2020}. Although these strategies maximize ridership by targeting major activity hubs, they may also contribute to persistent socio-demographic imbalances in the user base (see Section~\ref{sec:who_uses_bss})~\cite{Ogilvie2012, Woodcock2014, Ricci2015}.

These disparities are increasingly discussed within broader frameworks of transport equity and distributive justice~\cite{Cook2019, Litman2017, Pereira2017, Sheller2018}, and they can persist even in dockless systems. For example, in Shenzhen, China, both usage and perceived accessibility of dockless bicycles were higher in wealthier and more central neighborhoods~\cite{Zhou2024}. Reviews of cycling equity similarly report uneven distributions of infrastructure coverage, cycling activity, destination accessibility, and related health benefits, systematically favoring socially advantaged urban districts~\cite{Cunha2023}. Beyond physical access, informational barriers also matter: disadvantaged residents may be less aware of systems or targeted programs due to weaker user networks and limited access to smartphones or the internet~\cite{McNeil2018}. In Montreal, awareness of \textit{BIXI} increased mainly among higher-education groups and those living closer to stations, while lower-education groups remained less informed~\cite{Bernatchez2015}.

Encouragingly, evidence suggests that more equitable outcomes are achievable through inclusive design and policies. Interest in bike-sharing can be as high among disadvantaged groups as among advantaged ones~\cite{McNeil2018}, and per-capita usage can even be higher when barriers are reduced~\cite{Ogilvie2012}. In Washington, D.C., introducing \acp{e-b} alongside \acp{m-b} broadened the user base toward more racially diverse, lower-income, and more gender-balanced ridership~\cite{Bonilla-Alicea2020}. This implies that disparities are often driven less by preferences than by structural constraints, and that expanding coverage into underserved areas, maintaining affordability, and implementing targeted outreach can improve participation among marginalized communities~\cite{Goodman2014, McNeil2018, Ogilvie2012, Ursaki2015}.


 

    